\documentclass[11pt,letterpaper]{article}
\usepackage[margin=1in]{geometry}
\usepackage{amsmath, amssymb}
\usepackage{tikz}
\usepackage{enumitem}
\usepackage{hyperref}
\usepackage{listings}
\usepackage{xcolor}
\usepackage{longtable}
\setlength{\parindent}{0pt}
\setlength{\parskip}{1ex}

\title{Debian 13 Kiosk and Building Directory Deployment Guide}
\author{}
\date{\today}

\begin{document}

\maketitle

\section*{Purpose}

This document defines the standardized installation and configuration procedure for the Debian 13 (Trixie) kiosk display systems and building directory server environment.

It replaces all previous Debian 13 installation notes and integrates updated deployment practices including:

\begin{itemize}
\item Overlayroot read-only root filesystem
\item Autologin kiosk environment (no display manager)
\item Cage + Chromium Wayland kiosk
\item XFCE administrative fallback on USB keyboard insertion
\item Central directory server architecture
\item Removal of legacy monit usage
\end{itemize}

\section{System Architecture}

\subsection{Display Nodes}

Each display node runs:

\begin{itemize}
\item Debian 13 minimal install
\item Read-only overlayroot
\item Autologin kiosk user on tty1
\item Chromium kiosk pointing to central server
\item Optional XFCE admin mode triggered by USB keyboard
\end{itemize}

Example:

\begin{itemize}
\item Display node: 192.168.1.80
\item Server (dev): 192.168.1.131
\end{itemize}

\subsection{Directory Server}

NodeJS server providing:

\begin{itemize}
\item Static UI
\item REST API endpoints
\item Data storage (JSON or database)
\end{itemize}

Runs on port 3000.

\section{Debian 13 Base Installation}

\subsection{Install Debian 13}

Use official Debian 13 installer (amd64).

Select:

\begin{itemize}
\item Minimal system utilities only
\item SSH server
\item No desktop environment
\end{itemize}

Create user:

\begin{itemize}
\item username: kiosk
\end{itemize}

\subsection{Post-install packages}

\begin{lstlisting}[language=bash]
sudo apt update
sudo apt install \
xorg xfce4 xfce4-terminal \
chromium cage wlr-randr \
unclutter seatd \
overlayroot curl vim
\end{lstlisting}

\section{Overlayroot Configuration}

\subsection{Enable read-only root}

Edit:

\begin{lstlisting}[language=bash]
sudo vi /etc/overlayroot.conf
\end{lstlisting}

Set:

\begin{lstlisting}
overlayroot="tmpfs"
\end{lstlisting}

Reboot.

Verify:

\begin{lstlisting}[language=bash]
mount | grep overlayroot
\end{lstlisting}

\section{Autologin Configuration}

Disable display manager:

\begin{lstlisting}[language=bash]
sudo systemctl disable lightdm --now
sudo systemctl set-default multi-user.target
\end{lstlisting}

Configure tty1 autologin:

\begin{lstlisting}[language=bash]
sudo mkdir -p /etc/systemd/system/getty@tty1.service.d
sudo vi /etc/systemd/system/getty@tty1.service.d/override.conf
\end{lstlisting}

Insert:

\begin{lstlisting}
[Service]
ExecStart=
ExecStart=-/sbin/agetty --autologin kiosk --noclear %I $TERM
Type=simple
\end{lstlisting}

Reload:

\begin{lstlisting}[language=bash]
sudo systemctl daemon-reload
\end{lstlisting}

\section{Kiosk Startup Logic}

\subsection{.bash\_profile}

\begin{lstlisting}[language=bash]
if [[ "$(tty 2>/dev/null)" == "/dev/tty1" ]]; then
    export XDG_RUNTIME_DIR="/run/user/$(id -u)"
    mkdir -p "$XDG_RUNTIME_DIR"
    chmod 700 "$XDG_RUNTIME_DIR"

    while true; do
        rm -f /tmp/kiosk-exit
        /home/kiosk/building-directory/scripts/start-kiosk.sh

        if [[ -f /tmp/kiosk-exit ]]; then
            export XAUTHORITY=/tmp/.Xauthority
            export XDG_CONFIG_HOME=/tmp/xfce4-config
            export XDG_CACHE_HOME=/tmp/xfce4-cache
            export XDG_DATA_HOME=/tmp/xfce4-data
            mkdir -p /tmp/xfce4-config /tmp/xfce4-cache /tmp/xfce4-data

            startxfce4 -- -logfile /tmp/Xorg.0.log \
            >/tmp/xfce4-start.log 2>&1
        fi
    done
fi
\end{lstlisting}

\section{Cage + Chromium Startup}

\subsection{start-kiosk.sh}

\begin{lstlisting}[language=bash]
#!/bin/bash
SERVER_URL="http://192.168.1.131:3000"

cage -d -- sh -c '
OUTPUT=$(wlr-randr 2>/dev/null | sed -n "1s/ .*//p")
[ -n "$OUTPUT" ] && wlr-randr --output "$OUTPUT" --mode 1920x1080
rm -rf /tmp/chromium-profile 2>/dev/null || true

exec chromium \
--ozone-platform=wayland \
--user-data-dir=/tmp/chromium-profile \
--kiosk \
--noerrdialogs \
--disable-infobars \
--no-first-run \
--disable-session-crashed-bubble \
--overscroll-history-navigation=0 \
--disable-translate \
--touch-events=enabled \
'"$SERVER_URL"'
'
\end{lstlisting}

Make executable:

\begin{lstlisting}[language=bash]
chmod +x start-kiosk.sh
\end{lstlisting}

\section{USB Keyboard Admin Mode}

Udev rule detects keyboard insertion and triggers:

\begin{itemize}
\item Kill cage
\item Create /tmp/kiosk-exit
\item XFCE launches
\end{itemize}

On XFCE logout, kiosk restarts automatically.

\section{Server Deployment (NodeJS)}

\subsection{Directory server}

\begin{lstlisting}[language=bash]
cd ~/building-directory/server
npm install
node server.js
\end{lstlisting}

Verify:

\begin{lstlisting}[language=bash]
curl http://localhost:3000
\end{lstlisting}

\subsection{Allow LAN access}

If firewall enabled:

\begin{lstlisting}[language=bash]
sudo ufw allow 3000/tcp
\end{lstlisting}

\section{Testing Connectivity from Display}

\begin{lstlisting}[language=bash]
curl http://192.168.1.131:3000
curl http://192.168.1.131:3000/api/data-version
\end{lstlisting}

\section{Deployment Strategy}

\subsection{Recommended}

\begin{itemize}
\item Develop server on workstation (192.168.1.131)
\item Test kiosk connectivity
\item Freeze code
\item Deploy server to one display node or dedicated server
\end{itemize}

\subsection{Production Option}

Run server on one display node:

\begin{itemize}
\item Choose most reliable hardware
\item Assign static IP
\item Use systemd service for Node server
\end{itemize}

\section{Maintenance Checklist}

\begin{itemize}
\item Verify overlayroot active
\item Confirm kiosk autostarts after reboot
\item Confirm keyboard triggers admin XFCE
\item Confirm server reachable from kiosk
\item Confirm Chromium reloads after crash
\end{itemize}

\section{Notes}

\begin{itemize}
\item LightDM not used
\item Monit removed
\item Systemd handles all restart logic
\item Overlayroot protects filesystem
\end{itemize}

\end{document}
